\section{Итоги}

\begin{frame}[c]{Проверь себя}
  \begin{numrows}
    1 & Почему фото одного листа нельзя раскидать по train и test?\\
    2 & Кросс-валидация дала $0{,}84 \pm 0{,}03$ и $0{,}85 \pm 0{,}03$. Какая модель лучше?\\
    3 & $X$ имеет форму $500\times 4$. Какой формы $\nabla L$?\\
    4 & Чем шаг SGD отличается от шага ГД?\\
    5 & ГД расходится уже при маленьком шаге. Что проверить первым делом?\\
  \end{numrows}
\end{frame}

\begin{frame}[c]{Ответы}
  \begin{numrows}
    1 & Тест будет проверять память, а не обобщение: утечка\\
    2 & Сказать нельзя: разница меньше разброса\\
    3 & $4\times1$, как у $w$\\
    4 & Градиент по случайному батчу, а не по всем объектам\\
    5 & Масштаб признаков и формулу градиента\\
  \end{numrows}
\end{frame}

\begin{frame}[c]{На практике: восстановление коэффициентов}
  \small
  Даны точки $(x, f(x))$, где
  \[f(x) = \bigl((a+\varepsilon_a)\sin x + (b+\varepsilon_b)\ln x\bigr)^2 + (c+\varepsilon_c)\,x^2,
    \qquad \varepsilon \in [-0{,}001;\ 0{,}001].\]
  Нужно найти $a$, $b$, $c$ с точностью $10^{-2}$.

  \medskip
  \begin{columns}[T]
    \begin{column}{0.47\textwidth}
      \begin{block}{Подумайте}
        \small
        \begin{itemize}
          \item Что здесь данные, а что параметры?
          \item Чему равна $\partial f/\partial a$? А $\partial f/\partial c$?
        \end{itemize}
      \end{block}
    \end{column}
    \begin{column}{0.47\textwidth}
      \visible<2>{%
      \begin{block}{Подсказка}
        \small
        $\sin x$, $\ln x$, $x^2$ не зависят от $a$, $b$, $c$ ---
        это просто столбцы признаков:

        \smallskip
        \centering\footnotesize
        \zebra{2}\footnotesize
        \begin{tabular}{rrrr}
          $x$ & $\sin x$ & $\ln x$ & $x^2$\\
          $0{,}5$ & $0{,}479$ & $-0{,}693$ & $0{,}25$ \\
          $1$ & $0{,}841$ & $0{,}000$ & $1$ \\
          $2$ & $0{,}909$ & $0{,}693$ & $4$\\
        \end{tabular}
      \end{block}}
    \end{column}
  \end{columns}

  \smallskip
  {\footnotesize\color{acGray}Задача: \url{https://coderun.yandex.ru/problem/coefficients-restoration}\par}
\end{frame}
